\section{Seifert types, circle bundles and covers}\label{sec:seifert}

\subsection{Finite Seifert types}

Put $C=K(\Z,1)\simeq S^1$, with basepoint $0$, positive loop $h_C$ and multiplication maps $\mu_s:C\to C$ for $s\in\Z$. The pants graph is the HIT $P_0=C\hpo_1C$, with gluing path $\gamma$ and based loops $a=\operatorname{ap}_{\iota_1}(h_C)$, $b=\gamma\cdot\operatorname{ap}_{\iota_2}(h_C)\cdot\gamma^{-1}$. Let $\partial_1,\partial_2:C\to P_0$ be the inclusions and $\partial_3$ the map with loop $(ab)^{-1}$, specified by the fixed circle equivalence. Throughout, a \emph{marked equivalence of diagrams} retains the specified boundary maps, markings of core circles and bundle fibres, and the homotopies between them.


Set $T^2=C\times C$ and define $[s,n]:T^2\to C$ by $[s,n](u,t)=\mu_s(u)+\mu_n(t)$. The addition is that of the Eilenberg--Mac Lane model of $C$.
\begin{definition}\label{def:seifert}
For positive integers $p,q,r$, the finite Seifert type $M(p,q,r)$ is the pushout
\begin{equation}\label{diag:seifert}
\begin{tikzcd}[column sep=large,row sep=large]
\coprod_{i=1}^3 T^2 \arrow[r,"g"] \arrow[d,"f"'] \pomark & \coprod_{i=1}^3 C \arrow[d] \\
P_0\times C \arrow[r] & M(p,q,r),
\end{tikzcd}
\qquad
\begin{gathered}
f_i=(\partial_i,\id_C),\\
(g_1,g_2,g_3)=([1,p],[1,q],[1-r,r]).
\end{gathered}
\end{equation}
\end{definition}
The right vertices are solid-torus cores, with primitive kernel slopes $(p,-1)$, $(q,-1)$ and $(r,r-1)$. Van Kampen \cite[\S8.7]{hott} gives
\begin{equation}\label{eq:seifert-presentation}
\pi_1M(p,q,r)=\langle x,y,z,h\mid h\text{ central},\quad x^p=y^q=z^r=xyz=h\rangle,
\end{equation}
where $x,y,z$ are the filling cores and $h$ is the second-circle loop. The third boundary has image $z^{1-r}$ and $xy=z^{r-1}$.

\subsection{The two finite covers}

The four types and the finite groups used to construct the covers are
\[
\begin{array}{c@{\qquad}cc@{\qquad}cc@{\qquad}c}
 c&C_c&A_c&G_c&H_c&[G_c:H_c]\\ \hline
 O&M(2,3,4)&M(2,2,4)&\Sym_4&D_8&3\\
 I&M(2,3,5)&M(2,2,3)&\Alt_5&D_6&10.
\end{array}
\]
Here $G_c=\pi_1(C_c)/\langle h\rangle$. The subgroup $D_8$ is the normalizer of a four-cycle subgroup in $\Sym_4$, and $D_6$ is the normalizer of a three-cycle subgroup in $\Alt_5$; subscripts denote group orders. Fix these embeddings and finite coset representatives. The actions of the three boundary generators on $H_c\backslash G_c$ have orbit lengths
\begin{equation}\label{eq:orbit-lengths}
\begin{array}{c@{\qquad}ccc}
 &\partial_1&\partial_2&\partial_3\\ \hline
 O&2,1&3&1,2\\
 I&2,2,2,2,1,1&1,3,3,3&5,5.
\end{array}
\end{equation}

\begin{construction}[Coset covers of the Seifert diagram]\label{con:coset-cover}
For $H\leq G$, let $\mathcal Q_H(b)=\fib(\B H\to\B G,b)$ over $\B G=\mathrm{EM}_1G$, with coset fibre $H\backslash G$ and transport by right multiplication. A classifier $\rho:M(p,q,r)\to\B G$ gives the cover
\begin{equation}\label{eq:coset-total}
\widehat M_H=\sum_{m:M(p,q,r)}\mathcal Q_H(\rho m)\longrightarrow M(p,q,r).
\end{equation}
\end{construction}
We construct $\rho$ on \eqref{diag:seifert}. We use pushout flattening: the total space of a family over a pushout is the pushout of its restrictions \cite{hott,cubical}.

Let $(n_1,n_2,n_3)=(p,q,r)$ and $(s_1,s_2,s_3)=(1,1,1-r)$. Choose triangle generators $g_i\in G$ of exact orders $n_i$, with $g_1g_2g_3=1$. The auxiliary base is
\[
\mathcal O=P_0\hpo_{\coprod_i C}\coprod_i\B(\Z/n_i),
\]
with boundary maps $\partial_i$ and the cyclic-generator maps $C\to\B(\Z/n_i)$. Its map to $\B G$ sends $a,b$ to $g_1,g_2$ and the $i$th cyclic generator to $g_i$; the product relation supplies the third boundary comparison.

For a filling use the central extension
\[
A(n,s)=\Z^2/\langle(n,-s)\rangle
 =\langle v,h\mid [v,h]=1,\ v^n=h^s\rangle.
\]
The fibres of $\B A(n,s)\to\B(\Z/n)$ form an oriented circle family, trivialized on the boundary by the lifts $v,h$. For $\gcd(n,s)=1$, the isomorphism $A(n,s)\cong\Z$, $v\mapsto s$, $h\mapsto n$, gives boundary map $[s,n]$. The core maps to $s^{-1}\bmod n$, which is $1$ in all three fillings. Glue to the trivial family on $P_0$; univalence and flattening identify the total space with \eqref{diag:seifert}. Its composite through $\mathcal O\to\B G$ is $\rho$, with regular fibre sent to $1$ and specified boundary comparisons.

Flattening the pullback of $\mathcal Q_H$ expresses \eqref{eq:coset-total} as a pushout of lifted pieces. Over $\B\langle g_i\rangle$, components are indexed by $H\backslash G/\langle g_i\rangle$, with stabilizer
\[
S_t=\{k\in\langle g_i\rangle\mid tkt^{-1}\in H\},
\]
with its inclusion into $\langle g_i\rangle$ and conjugation map to $H$. An orbit of length $\ell$ has stabilizer of order $d=n_i/\ell$, generated by $g_i^\ell$. The extension restricts to $A(d,s_i)$, by $v'\mapsto v^\ell$, $h'\mapsto h$. Circle-orbit charts therefore give the actual boundary map
\begin{equation}\label{eq:boundary-cover}
T^2\longrightarrow T^2,\qquad (u,t)\longmapsto(\mu_\ell(u),t),
\end{equation}
including its comparison with both adjacent vertices.

The lifted diagrams give the following result.
\begin{theorem}[Short covers]\label{thm:short-cover}
There are specified equivalences
\[
\theta_c:A_c\simeq\sum_{m:C_c}\mathcal Q_{H_c}(\rho_c m).
\]
The maps $p_c=\operatorname{pr}_1\theta_c:A_c\to C_c$ have set fibres of cardinality $N_O=3$ and $N_I=10$.
\end{theorem}
\begin{proof}
\emph{Cancellation and retained boundaries.}
A regular filling has multiplicity $1$, so its base vertex
$\B(\Z/1)$ is contractible. After contracting a spanning tree, a filling
whose boundary contains a graph generator $t$ exactly once has word
$u t^{\epsilon}v$, with $\epsilon\in\{1,-1\}$ and $u,v$ in the remaining graph.
The elementary cancellation removes this loop together with its filling cell.
Indeed, for any target $Z$, after fixing a map from the remaining graph,
the extension data are a loop $\ell$ and a path
$u'\ell^{\epsilon}v'=\mathsf{refl}$, where $u',v'$ are the images of $u,v$.
The function $\ell\mapsto u'\ell^{\epsilon}v'$ is an equivalence of the loop type,
so these data are contractible. The pushout mapping property therefore makes
the inclusion of the remaining graph an equivalence.

Reindexing the circle family along this base equivalence gives the
corresponding total-space equivalence. On the remaining graph its restriction
is the specified trivial family. Write $i$ for this total-space equivalence
and $r$ for its inverse. A retained boundary map $b$ is carried to $r b$,
with comparison $i r b\simeq b$ inherited from the inverse law. Thus the comparison includes the whole
torus map. The changes of fibre coordinates are the shears computed below.
These are the admissibility conditions checked for the primitive cancellations
in the finite tables.

We give the finite reduction explicitly. Permutations below compose as functions, with the right factor applied first. The triangle generators are
\[
\begin{array}{c@{\qquad}ccc}
 &g_1&g_2&g_3=(g_1g_2)^{-1}\\ \hline
 O&(12)&(234)&(1432)\\
 I&(12)(34)&(235)&(15342).
\end{array}
\]
In case $O$, number the sheets $0,1,2$ using representatives $1,g_1,g_2$. The three boundary orbit lists are $(0,1),(2)$; $(0,2,1)$; and $(0),(1,2)$. Contract the tree formed by the $a$- and $b$-edges from sheet $0$. The four remaining graph generators, written as words in the original pants loops, are
\[
\xi_1=a^2,\qquad \xi_2=ab,\qquad
\xi_3=bab^{-1},\qquad \xi_4=b^2a^{-1}.
\]
For each boundary write $(w,k;[s,d])$ for the left torus loops $(wh^k,h)$ and the right row $[s,d]$; its filling relation is $(wh^k)^d=h^s$. A label $A_i,B_i,C_i$ means the orbit starting at sheet $i$ for $\partial_1,\partial_2,\partial_3$, respectively. The initial five fillings are
\[
\begin{array}{c@{\quad}c@{\quad}r@{\quad}c}
 &w&k&[s,d]\\ \hline
 A_0&\xi_1&0&[1,1]\\
 A_2&\xi_3&0&[1,2]\\
 B_0&\xi_4\xi_2&0&[1,1]\\
 C_0&\xi_2^{-1}&0&[-3,4]\\
 C_1&\xi_4^{-1}\xi_3^{-1}\xi_1^{-1}&0&[-3,2].
\end{array}
\]
There are exactly two primitive cancellations: $A_0$ eliminates $\xi_1$ by $\xi_1=h$, and $B_0$ eliminates $\xi_2$ by $\xi_2=\xi_4^{-1}h$. With $u=\xi_3$ and $v=\xi_4$, the surviving data are
\[
\begin{array}{c@{\qquad}ccc}
 &A_2&C_0&C_1\\ \hline
 w&u&v&v^{-1}u^{-1}\\
 k&0&-1&-1\\
 {[s,d]}&[1,2]&[-3,4]&[-3,2].
\end{array}
\]
Choose the free basis $U=v^{-1}u^{-1}$, $V=u$, whose inverse is $u=V$, $v=V^{-1}U^{-1}$, and order the boundaries as $(C_1,A_2,C_0)$. Their words become $U,V,(UV)^{-1}$. Precomposing each torus by $(a,t)\mapsto(a,t-\mu_k(a))$ removes its left winding $k$ and changes the right weight to $s-dk$. Writing $\nu_i$ for the ordered filling multiplicities, we obtain
\[
\begin{aligned}
(\nu_1,\nu_2,\nu_3)&=(2,2,4),
& (k_1,k_2,k_3)&=(-1,0,-1),\\
(s_1,s_2,s_3)&=(-3,1,-3),
& (s_i-\nu_ik_i)_{i=1}^3&=(-1,1,1).
\end{aligned}
\]
Finally change the graph trivialization by winding $1$ on $U$ and $0$ on $V$. The induced boundary shear integers are $(1,0,-1)$, so the weights become $(1,1,-3)$. In the old total-space coordinates the final pants generators are
\[
A=Uh=v^{-1}u^{-1}h,\qquad B=V=u,\qquad
(AB)^{-1}=vh^{-1}.
\]
This gives precisely the rows $[1,2],[1,2],[-3,4]$ of $M(2,2,4)$.

For case $I$, the following table records the corresponding finite data. Sheets are numbered $0,\ldots,9$; the graph letters $\xi_i$ are indexed by the non-tree edges in increasing sheet order, with the $a$-edge before the $b$-edge. The chosen tree consists of the $a$-edges at sheets $0,2,3$ and the $b$-edges at sheets $1,3,4,5,7,9$.
\begin{center}
\small
\begin{tabular}{@{}ll@{}}
\toprule
Data & Case $I$ \\
\midrule
$\partial_1$ orbits & $(0,1),(2,4),(3,5),(6,9),(7),(8)$ \\
$\partial_2$ orbits & $(0),(1,2,3),(4,6,7),(5,8,9)$ \\
$\partial_3$ orbits & $(0,1,5,6,2),(3,4,7,9,8)$ \\
Initial graph rank; fillings & $11;\ 12$ \\
Cancelled boundaries & $A_0,A_2,A_3,B_1,B_4,B_5,A_6,C_0,C_3$ \\
Eliminated generators, in order & $\xi_2,\xi_4,\xi_5,\xi_3,\xi_7,\xi_{10},\xi_6,\xi_1,\xi_8$ \\
Surviving generators & $u=\xi_9,\ v=\xi_{11}$ \\
Surviving boundaries & $A_7:(v^{-1}u^{-1},0;[1,2])$ \\
 & $A_8:(u,0;[1,2]),\quad B_0:(v,1;[1,3])$ \\
Final basis and inverse & $U=v^{-1}u^{-1},\ V=u;\quad u=V,\ v=V^{-1}U^{-1}$ \\
Left windings; weights after removal & $(0,0,1);\quad(1,1,-2)$ \\
Final graph gauge & $(0,0)$ \\
\bottomrule
\end{tabular}
\end{center}
There are exactly nine primitive cancellations. The same basis substitution and the torus shears $(0,0,-1)$ give $[1,2],[1,2],[-2,3]$, the rows of $M(2,2,3)$. The graph ranks $4,11$ also follow from $2N_c$ edges and $N_c$ vertices, and the orbit lists give the filling counts $5,12$.

These are equivalences of the untruncated pushouts. The routine moves are tree contraction, cancellation of an identity leg or a primitive loop--cell pair, graph maps with the displayed inverse substitutions, and torus shears. The finite certificates in Appendix~\ref{app:formal} supply the orbit and tree data, the boundary words, the legality of each cancellation, and the terminal integer calculations; their interpretation retains the specified boundary homotopies and torus squares. In particular, the surviving core circles are unchanged, apart from their reordering: $(C_1,A_2,C_0)$ in case $O$ and $(A_7,A_8,B_0)$ in case $I$. The regular fibre $h$ is retained throughout. Composing the diagram equivalences gives
\[
\theta_c:A_c\simeq\sum_{m:C_c}\mathcal Q_{H_c}(\rho_c m).
\]
The coset families have constant finite cardinality, so their projections give the asserted covers.
\end{proof}

\subsection{The twisted circle bundle}

Write $e=4$ in case $O$ and $e=3$ in case $I$. Staging the first two fillings of $M(2,2,e)$ and cancelling their identity legs gives
\begin{equation}\label{eq:staged-twisted}
W=C_x\hpo_{C_h}C_y\quad(h\mapsto x^2,\ h\mapsto y^2),
\qquad
M(2,2,e)\simeq W\hpo_{T^2}C,
\end{equation}
whose remaining boundary has loops $((xy)^{-1},h)$ and right leg $[1-e,e]$. Let $D_{\mathrm{inv}}:C\to\mathcal U$ have fibre $C$ and transport $t\mapsto-t$ on the positive loop. Its total space $K_{\mathrm{inv}}=\sum_{u:C}D_{\mathrm{inv}}(u)$ is the inversion mapping torus. The Nielsen change
\[
x\mapsto j,\qquad y\mapsto jr^{-1},\qquad
j\mapsto x,\quad r\mapsto y^{-1}x
\]
identifies $W\simeq K_{\mathrm{inv}}$: circle and pushout elimination identify the relations $x^2=y^2$ and $jr=r^{-1}j$, including the relation square and inverse homotopy. Now $h=j^2$, $(xy)^{-1}=rj^{-2}$, and the torus equivalence
\[
\sigma(u,v)=(v,u+v),\qquad \sigma^{-1}(a,t)=(t-a,a)
\]
changes the left boundary to $(j^2,r)$ and the right row to $[e,1]$, since $(1-e)v+e(u+v)=eu+v$. Hence
\begin{equation}\label{eq:twisted-normal-form}
M(2,2,e)\simeq K_{\mathrm{inv}}\hpo_{T^2}C,
\qquad T^2\xrightarrow{(j^2,r)}K_{\mathrm{inv}},\quad T^2\xrightarrow{[e,1]}C.
\end{equation}

Let $T:J\to\mathcal U$ be the orientation two-point family on $J=\RP^1$ and $\widetilde J=\sum_{x:J}T(x)$. Fix circle charts with $\widetilde J\to J$ of degree $2$ and $f_A:\widetilde J\simeq C$. The projective-space construction gives $R=\RP^2=J\hpo_{\widetilde J}1$ \cite{cubical}. For a two-element type $T$ with involution $z\mapsto\bar z$, put
\[
\Z_T=\{v:T\to\Z\mid v(\bar z)=-v(z)\},\qquad
\operatorname{SignC}(T)=K(\Z_T,1).
\]
Evaluation at $z:T$ gives $\epsilon_z:\operatorname{SignC}(T)\simeq C$, and $\epsilon_{\bar z}=-\epsilon_z$. Define $\mathcal F_e:R\to\mathcal U$ by $\operatorname{SignC}(T(x))$ on $J$, by $C$ on the cap, and by the seam equivalence
\begin{equation}\label{eq:sign-clutching}
v\longmapsto\epsilon_z(v)+\mu_e(f_A(x,z)),\qquad (x,z):\widetilde J.
\end{equation}
Univalence defines the family from this clutching; its dependence on $T$ supplies the deck coherences. Flattening gives
\[
M_e:=\sum_{b:R}\mathcal F_e(b)\simeq K_{\mathrm{inv}}\hpo_{T^2}C
\]
with exactly the legs in \eqref{eq:twisted-normal-form}. This defines $q_e:M_e\to R$ and, by the specified equivalence, $q_c:A_c\to R$.

The fibre at the cap is $C$, with generator $r$, and transport around the nontrivial base loop is inversion. The source group and bundle character are
\begin{equation}\label{eq:bundle-character}
\pi_1(M_e)=\langle r,j\mid jrj^{-1}=r^{-1},\ j^2=r^e\rangle,
\qquad r\mapsto0,\quad j\mapsto1\quad\text{in }\Z/2.
\end{equation}
The third filling core is unchanged in \eqref{eq:twisted-normal-form}. If $x,y,z$ are the positive loops of the three source cores, the comparison therefore gives
\[
x\longmapsto j,\qquad y\longmapsto jr^{-1},\qquad
z\longmapsto r,\qquad h\longmapsto j^2=r^e.
\]
Thus the bundle fibre is marked by $z$, and the regular Seifert fibre is $h=r^e$.

\begin{proposition}\label{prop:source-pi2}
For $e=3,4$, the circle bundle $M_e\to R$ has $\pi_2(M_e)=0$.
\end{proposition}
\begin{proof}
The proof has three steps: pull the bundle back to $S^2$; construct its classifying map with nonzero clutching winding; and use injectivity on $\pi_2$ and the fibre sequence to deduce vanishing.

\emph{Pullback to the sphere.} Let
\[
\widetilde M_e=\sum_{s:S^2}\mathcal F_e(qs),
\qquad Q:\widetilde M_e\longrightarrow M_e,
\qquad Q(s,v)=(qs,v),
\]
where $q:S^2\to R$ is the orientation cover. The square formed by $Q$, $q$ and the two bundle projections is cartesian. We first compute $\pi_2(\widetilde M_e)$.

In the two-cap chart for this cover, the base is
$\widetilde J\hpo_{\widetilde J\times\{0,1\}}\{0,1\}\simeq\Sigma C$.
Evaluation in the chosen sheet identifies the circle fibres on the equator and the caps with $C$. The two seam equivalences are translations by
\[
s_0(a)=\mu_e(f_A(a)),\qquad
s_1(a)=\mu_{-e}(f_A(\delta a)),
\]
where $\delta$ is the deck map of $\widetilde J\to J$. The circle chart has a specified homotopy $f_A\delta\simeq f_A$, including its vertex paths. It identifies these seams with translations by $\mu_e$ and $\mu_{-e}$ in the common equatorial coordinate (Figure~\ref{fig:two-cap}). In the chosen suspension coordinate the clutching winding is $-2e$; only its nonvanishing is used.

\begin{figure}[htbp]
\centering
\begin{tikzcd}[column sep=huge,row sep=large]
C_0 \arrow[d] & C_{\mathrm{eq}}\times C
  \arrow[l,"{(a,t)\mapsto t+\mu_e(a)}"']
  \arrow[r,"{(a,t)\mapsto t+\mu_{-e}(a)}"]
  \arrow[d,"\operatorname{pr}_1"] & C_1 \arrow[d] \\
1 & C_{\mathrm{eq}} \arrow[l] \arrow[r] & 1
\end{tikzcd}
\caption{The pulled-back circle bundle in two-cap coordinates. The lower homotopy pushout is $\Sigma C\simeq S^2$ and the upper one is $\widetilde M_e$. The seam translations have windings $+e,-e$; cancelling the first cap leaves their difference $(-e)-e=-2e$.}
\label{fig:two-cap}
\end{figure}

\emph{Classifying map and winding.} Put $K=K(\Z,2)$ and let $\omega:C\simeq\Omega(K,0)$ be the pointed circle equivalence, with its specified basepoint comparison. For the path family $U(k)=(0=k)$, transport along a loop $\omega(s)$ is right concatenation by $\omega(s)$. The compatibility of $\omega$ with addition therefore identifies translation by $s$ on $C$ with this transport. Assigning $0$ to the equator and both caps, and $\omega(s_b(a))$ to the $b$th seam, defines a map to $K$ whose pulled-back path family is the circle family above. On cancelling the first cap, the remaining seam is
\[
\omega(\mu_e(a))^{-1}\cdot\omega(\mu_{-e}(a))
 =\omega(\mu_{-2e}(a)),
\]
where the equality is the homotopy supplied by the addition and inverse laws. Thus, with $d=-2e$, the resulting map is
\[
c:\Sigma C\longrightarrow K,
\qquad c(\mathsf{north})=c(\mathsf{south})=0,
\qquad \operatorname{ap}_c(\mathsf{merid}(a))=\omega(\mu_d(a)).
\]
The marked cancellation retains the seam homotopies, so it gives an equivalence of total spaces
\[
\widetilde M_e\simeq\sum_{a:\Sigma C}(0=c(a))
 \simeq\fib(c,0)=:V_e.
\]
The second equivalence reverses the path in each fibre.

\emph{Injectivity and the fibre sequence.} To establish injectivity of $c_*$ on $\pi_2$, let
$\sigma:C\to\Omega(\Sigma C,\mathsf{north})$ be the suspension unit,
$\sigma(a)=\mathsf{merid}(a)\cdot\mathsf{merid}(0)^{-1}$,
with its basepoint comparison, and put
$w=\sigma_*[h_C]\in\pi_1\Omega\Sigma C=\pi_2(\Sigma C)$.
The suspension--loop adjunction gives the pointed homotopy
$\Omega c\circ\sigma\simeq\omega\circ\mu_d$.
Choose the integer coordinate on $\pi_2K$ induced by $\omega$ and the positive generator of $\pi_1C$. Since $(\mu_d)_*[h_C]=d[h_C]$, this homotopy gives
\[
c_*(w)=d=-2e\neq0
\]
in that target coordinate. Also $\Sigma C\simeq S^2$, so $\pi_2(\Sigma C)\cong\Z$, and $\pi_2K\cong\Z$. Under any group isomorphisms with $\Z$, the homomorphism $c_*$ has the form $n\mapsto kn$. Its nonzero value on $w$ forces $k\neq0$, and cancellation in $\Z$ then makes it injective. Only a nonzero image is used; the sign and orientation of this suspension coordinate are not compared with those of the standard sphere.

The fibre sequence for $c$ gives
\[
0=\pi_3K\longrightarrow\pi_2(V_e)
 \longrightarrow\pi_2(\Sigma C)
 \xrightarrow{c_*}\pi_2K.
\]
Exactness and injectivity of $c_*$ give $\pi_2(V_e)=0$. The fibre $V_e$ is connected because $\Sigma C$ is connected and $\pi_1K=0$, so vanishing transfers to $\widetilde M_e$ at every basepoint. Finally $\fib(Q,(b,v))\simeq\fib(q,b)$ is a two-element set; hence $Q$ induces an isomorphism on $\pi_2$. Thus $\pi_2(M_e)=0$ for $e=3,4$, with basepoint transport used only in the proposition expressing vanishing.
\end{proof}

\subsection{Truncation squares and the common fibre}\label{sec:truncation}

\begin{theorem}[Truncation squares]\label{thm:trunc-square}
Let $n\geq0$ and $k=n+1$. If every fibre of $f:X\to Y$ is an $n$-type, its naturality square with $\eta_k$ is cartesian. Explicitly, for the canonical path $c_x:\tau_kf(\eta_kx)=\eta_k(fx)$, the comparison
\begin{equation}\label{eq:trunc-comparison}
X\longrightarrow\sum_{y:Y}\sum_{a:\tau_kX}(\tau_kf(a)=\eta_ky),
\qquad x\longmapsto(fx,\eta_kx,c_x)
\end{equation}
is an equivalence.
\end{theorem}
\begin{proof}
The universe $\mathcal U^{\leq n}$ is an $(n+1)$-type \cite[\S7.1]{hott}, so the classifying map $y\mapsto\fib(f,y)$ of $f$ extends along $\eta_k:Y\to\tau_kY$ to a family $D$ over $\tau_kY$ with $D(\eta_ky)\simeq\fib(f,y)$. The total space $\sum_tD(t)$ is a $k$-type, and truncation induction on $\tau_kX$ and on $\tau_kY$ gives mutually inverse maps $\tau_kX\rightleftarrows\sum_tD(t)$ over $\tau_kY$. The pullback of $\sum_tD(t)\to\tau_kY$ along $\eta_k$ is $\sum_y\fib(f,y)\simeq X$, which gives \eqref{eq:trunc-comparison}.
\end{proof}
Taking fibres in the other direction gives
\begin{equation}\label{eq:unit-fibres}
\fib(\eta_k:X\to\tau_kX,a)\simeq
\fib(\eta_k:Y\to\tau_kY,\tau_kf(a)).
\end{equation}

\begin{lemma}[Pointed Hopf reflection]\label{lem:hopf-reflection}
The pointed Hopf classifying map $h:S^2\to_*K(\Z,2)$ induces a pointed equivalence
\[
\bar h:\tau_2S^2\eqp K(\Z,2),\qquad \bar h\circ\eta_2\sim_*h.
\]
There is an unpointed equivalence
$\fib(\eta_2:S^2\to\tau_2S^2,\eta_2(*))\simeq S^3$.
\end{lemma}
\begin{proof}
The Hopf construction gives $h$ with basepoint fibre $S^3$ \cite[\S8.5]{hott}. Since $K(\Z,2)$ is connected, the propositional assertion that its fibres are $2$-connected extends from the basepoint to every point. Thus $h$ is $2$-connected with a $2$-truncated target, so it exhibits the $2$-truncation of $S^2$. The universal property gives $\bar h$ and its unit comparison; the pointing of $h$ supplies the pointing of $\bar h$ and compatibility of the comparison. Taking the basepoint fibre in this square gives the stated sphere equivalence.
\end{proof}
Applying \eqref{eq:unit-fibres} to the orientation cover $q:S^2\to R$ at level $2$ now gives:
\begin{proposition}[Common fibre]\label{prop:common-F}
For $r_0=q(*)$ and $\baseN=\tau_2R$, there is an equivalence
$F=\fib(\eta_2:R\to\baseN,\eta_2r_0)\simeq S^3$.
\end{proposition}
The base $\baseN$ has $\pi_1=\Z/2$ and $\pi_2=\Z$ with sign action.

For $q_c$ use Theorem~\ref{thm:trunc-square} at level $2$ and the equivalence $\tau_2A_c\simeq\tau_1A_c$ supplied by Proposition~\ref{prop:source-pi2}. For $p_c$ use level $1$. This gives the two mixed cartesian naturality squares on the left of \eqref{diag:mixed-squares}; $\bar q_c$ denotes the induced map to $\baseN$.

\begin{proposition}[Cartesian pushout descent]\label{prop:pushout-descent}
A map of spans whose two naturality squares are cartesian induces cartesian squares from each vertex to the corresponding pushouts.
\end{proposition}
The cartesian squares identify the source span with the total spaces of compatible families over the target span. Glue these families by pushout induction; the flattening lemma identifies their total space with the source pushout, and its restrictions to the vertices give the asserted cartesian squares \cite[\S6.12]{hott}.

Apply descent to $R\xleftarrow{q_c}A_c\xrightarrow{p_c}C_c$ and $\baseN\xleftarrow{\bar q_c}\tau_1A_c\xrightarrow{\tau_1p_c}\tau_1C_c$. Their pushouts are $E_c$ and $P_c$, as in (S2), and the third square below is cartesian:
\begin{equation}\label{diag:mixed-squares}
\begin{tikzcd}[column sep=1.4em,row sep=large]
A_c \arrow[r,"q_c"] \arrow[d,"\eta_1"'] \pbmark & R \arrow[d,"\eta_2"]
&[-.3em] A_c \arrow[r,"p_c"] \arrow[d,"\eta_1"'] \pbmark & C_c \arrow[d,"\eta_1"]
&[-.3em] R \arrow[r] \arrow[d,"\eta_2"'] \pbmark & E_c \arrow[d,"t_c"] \\
\tau_1A_c \arrow[r,"\bar q_c"'] & \baseN
& \tau_1A_c \arrow[r,"\tau_1p_c"'] & \tau_1C_c
& \baseN \arrow[r] & P_c.
\end{tikzcd}
\end{equation}
It identifies the specified fibre of $t_c$ with $F$ and gives fibre sequences $F\to E_c\xrightarrow{t_c}P_c$. This proves (S1) and (S2).

\begin{lemma}[Dihedral-character rigidity]\label{lem:dihedral-character}
Let $Q$ be a group generated by $r,j$ with $jrj^{-1}=r^{-1}$, and let $\chi:Q\to\Z/2$ be a homomorphism satisfying $\chi(r)=0$ and $\chi(j)=1$. For $m\geq3$, let $\epsilon:D_{2m}\to\Z/2$ be the rotation/reflection character. If a homomorphism $\lambda:Q\to D_{2m}$ has nonabelian image, then $\chi=\epsilon\circ\lambda$.
\end{lemma}
\begin{proof}
Put $a=\lambda(r)$ and $b=\lambda(j)$. They do not commute, since they generate the image. The relation $bab^{-1}=a^{-1}$ implies that $a^2=1$ would force commutation. Every reflection squares to $1$, so $a$ is a rotation. Since rotations commute, $b$ is a reflection. The two characters agree on the generators.
\end{proof}

\begin{proposition}\label{prop:seifert-simple}
The types $E_c$ and $P_c$ are simply connected.
\end{proposition}
\begin{proof}
The types $R$, $A_c$ and $C_c$ are connected. For $R$ this follows from its projective-space pushout; for the Seifert types, the connected vertex $P_0\times C$ is joined to each connected filling core by a nonempty boundary torus. Choose a basepoint $a_0:A_c$, use its images as basepoints in $R$ and $C_c$, and use the pushout seam at $a_0$ to compare their images in $E_c$. Van Kampen \cite[\S8.7]{hott} identifies $\pi_1(E_c)$ with the pushout of groups
\[
\pi_1(R)\xleftarrow{\chi_c}\pi_1(A_c)
 \xrightarrow{(p_c)_*}\pi_1(C_c),
\qquad \chi_c=(q_c)_*:\pi_1(A_c)\longrightarrow\Z/2.
\]
The bundle character is surjective, since \eqref{eq:bundle-character} gives $\chi_c(j)=1$. Equivalently, its surjectivity follows from the connected circle fibres of $q_c$ and the low-degree homotopy exact sequence.

Write $L_c=\pi_1(C_c)$ and let
\[
K_c=\bigl\langle\!\bigl\langle
 (p_c)_*(\ker\chi_c)
\bigr\rangle\!\bigr\rangle_{L_c}
\]
be the normal closure in $L_c$. Since $\chi_c$ is surjective, a compatible pair of homomorphisms from the span to a group $L$ is determined by a homomorphism $v:L_c\to L$ for which $v\circ(p_c)_*$ kills $\ker\chi_c$. Equivalently, $v$ kills $K_c$. The two universal properties therefore give
\[
\pi_1(E_c)\cong L_c/K_c.
\]

We identify elements of $K_c$ using the actual core lifts. The coset-cover model has a map to $\B H_c$, inducing
$\lambda_c:\pi_1(A_c)\to H_c$ after transport along $\theta_c$. The marked classifier sends the third target core in case $O$ to the chosen four-cycle rotation in $H_O=D_8$, and the second target core in case $I$ to the chosen three-cycle rotation in $H_I=D_6$. Their based lifts give elements $u_c\in\pi_1(A_c)$ satisfying
\[
(p_O)_*(u_O)=z,\qquad (p_I)_*(u_I)=y,
\qquad \lambda_c(u_c)=\rho_c^{\mathrm{rot}},
\]
where $\rho_c^{\mathrm{rot}}$ denotes the chosen rotation in $H_c$. Indeed, in the pullback of $\B H_c\to\B G_c$, a target loop whose classifier is the image of a specified element of $H_c$ lifts using that element and the classifier comparison. The same construction lifts the marked reflection word, so the image of $\lambda_c$ contains the chosen rotation and reflection. These two elements do not commute in either $D_8$ or $D_6$.

Let $\epsilon_c:H_c\to\Z/2$ be the dihedral character, equal to $0$ on rotations and $1$ on reflections. We claim
\[
\chi_c=\epsilon_c\circ\lambda_c.
\]
Apply Lemma~\ref{lem:dihedral-character} to the generators supplied by the marked bundle comparison \eqref{eq:bundle-character} and the noncommuting rotation and reflection in the image of $\lambda_c$; their mere existence suffices because equality of homomorphisms is a proposition.
In particular,
$\chi_c(u_c)=\epsilon_c(\rho_c^{\mathrm{rot}})=0$.
Hence $z\in K_O$ and $y\in K_I$.

It remains to compute the quotients using \eqref{eq:seifert-presentation}; denote the images of its generators by the same letters. In both cases $x^2=xyz$ implies $x=yz$, and therefore $z=y^{-1}x$.
For $O$, the relation $z=1$ gives $x=y$. Then
\[
y^2=x^2=y^3
\]
forces $y=1$ by cancellation, whence $x=1$ and $h=x^2=1$. Thus all four generators of $L_O/K_O$ are trivial.
For $I$, the relation $y=1$ gives $h=y^3=1$ and $x^2=h=1$. Also $z=y^{-1}x=x$, so $z^5=h$ gives $x^5=1$. Since
\[
x^5=(x^2)^2x=x,
\]
we obtain $x=z=1$. Again all generators are trivial. Consequently $\pi_1(E_c)=1$ in both cases.

The pushout $E_c$ is connected because its vertices are connected and its overlap is inhabited. The vertices $\baseN=\tau_2R$ and $\tau_1C_c$ of $P_c$ are also connected, and $\tau_1A_c$ is inhabited, so $P_c$ is connected. Finally the fibre sequence $F\to E_c\to P_c$ gives the exact segment
\[
\pi_1(E_c)\longrightarrow\pi_1(P_c)\longrightarrow\pi_0(F).
\]
Here $\pi_1(E_c)=1$ and $\pi_0(F)=1$, since $F\simeq S^3$ is connected. Exactness therefore gives $\pi_1(P_c)=1$. Connectedness and triviality of these fundamental groups prove that $E_c$ and $P_c$ are simply connected.
\end{proof}
